\chapter{EXPECTED RESULTS}

\section{Deliverables}

\begin{enumerate}[leftmargin=*]
    \item \textbf{A longitudinal Vietnamese hotel rate dataset.} At least 100,000 price observations covering properties in five cities, each row carrying both an observation timestamp and a stay date, with room and rate-plan identity hashes, reference match status, availability status and provenance fields. The pilot averaged 9.2 stored observations per successful item at a 92.6\% success rate once the one unbookable property is removed from the cohort, so a full-cohort crawl day across 355 properties and 5 check-in dates produces on the order of 15,000 observations. On that basis the 100,000 target is reached within roughly seven full-cohort crawl days, and collection continues well past it to lengthen the series rather than only to enlarge it. The dataset ships with a data dictionary and a data-quality report covering missingness, anomaly rate and reference coverage per city.

    \item \textbf{Four trained and compared models.} XGBoost, Random Forest, LSTM and Transformer, each trained for the 1, 3, 7 and 14-day horizons, with hyperparameter configurations, training curves and held-out test metrics. Every result is reproducible from a configuration file and a fixed seed.

    \item \textbf{A comparative analysis report.} Accuracy per model and horizon, breakdown by city, review-score tier and lead-time bucket, SHAP feature attribution for the selected model, and an ablation study measuring the contribution of the calendar, lead-time, price-history, competitive-set and external-context groups.

    \item \textbf{A working system.} The internal operations application (already implemented) plus the product dashboard: a FastAPI backend serving both route groups, a Python crawl worker over the MySQL durable queue, two Next.js applications, and a Docker Compose configuration for deployment.

    \item \textbf{Two decision simulations.} A pricing opportunity analysis quantifying detected mispricing for the hotelier view, and a booking savings simulation quantifying the cost difference from following the timing signal for the consumer view. Both run retrospectively on the held-out period.

    \item \textbf{The thesis and defence materials.}
\end{enumerate}

\section{Performance Targets}

Table~\ref{tab:expected_performance} lists the target for each metric together with the reasoning that fixes it at that value.

\begin{table}[!htbp]
\centering
\caption{Target values and their basis}
\label{tab:expected_performance}
\begin{tabular}{|p{4.4cm}|p{2.6cm}|p{6.6cm}|}
\hline
\textbf{Metric} & \textbf{Target} & \textbf{Basis} \\
\hline
MAPE, 7-day horizon & $\leq 12\%$ & Within the range reported for comparable hotel price and tourism demand forecasting work \cite{law2019, zhang2011}. Primary metric because it is scale-invariant across properties whose rates differ by an order of magnitude. \\
\hline
RMSE, 7-day horizon & $\leq 350{,}000$ VND & Roughly 10\% of a typical nightly rate in the 3 to 5-star segment covered here. \\
\hline
$R^2$ & $\geq 0.75$ & The model accounts for at least three quarters of rate variance on held-out data. \\
\hline
Price direction accuracy & $\geq 80\%$ & Direction is the signal both views act on. Four correct calls in five is the threshold at which acting on the signal beats ignoring it. \\
\hline
Simulated booking savings & $\geq 15\%$ & Measured against booking at the first observed price on the same stay date. \\
\hline
Crawl success rate & $\geq 95\%$ & Per run, counting technical failures only. Sold-out nights and unbookable properties are valid outcomes, not failures. \\
\hline
API latency, p95 & $\leq 200$ ms & Dashboard responsiveness under simulated concurrent load. \\
\hline
Data freshness & $\leq 1$ hour & Time between a crawl completing and its observations being visible in the dashboard. \\
\hline
Retraining time & $\leq 30$ minutes & Full retraining on the accumulated dataset. \\
\hline
\end{tabular}
\end{table}

The accuracy targets are hypotheses to be tested, not claims. If a model family falls short, the result is reported as measured and the gap analysed by horizon, lead-time bucket and city, since a model that forecasts well at 14 days and poorly at 1 day, or well in Hanoi and poorly in Phu Quoc, is a more useful finding than a single averaged number.

\section{Contributions}

\begin{itemize}[leftmargin=*]
    \item A longitudinal Vietnamese hotel rate dataset with two independent time coordinates, which no public source currently provides.
    \item A room-comparability mechanism that keeps a rate series describing one product over months without manual per-property configuration, together with the evidence that a naive deduplication rule fails at this task.
    \item A comparison of four model families on short-horizon hotel rate forecasting for a Southeast Asian market, with feature-group attribution.
    \item A rate-intelligence design aimed at independent operators rather than at chains, evaluated by simulation on held-out data rather than by feature listing.
\end{itemize}
